% Authoritative initial attention probabilities from teaching_reference.py.
\tagstructbegin{tag=Figure,alttext={三个位置的因果注意力热图。第一行只有 Key 1 权重 1，其余未来位置被掩码；第二行的 Key 1 和 Key 2 权重为零点一二四四和零点八七五六；第三行权重为零点一六八零、零点二四一八和零点五九零三。}}
\begin{center}
\begin{tikzpicture}[x=15mm,y=12mm]
\foreach \r in {0,1,2}{
 \node[font=\footnotesize,anchor=east] at (-.12,{2.5-\r}) {Query \number\numexpr\r+1\relax};
}
\foreach \c in {0,1,2}{
 \node[font=\footnotesize,anchor=south] at ({\c+.5},3.06) {Key \number\numexpr\c+1\relax};
}
\filldraw[fill=datacolor!85,draw=linegray] (0,2) rectangle (1,3);
\node[white,font=\small\bfseries] at (.5,2.5) {$1.000$};
\filldraw[fill=paper,draw=linegray] (1,2) rectangle (2,3);
\node[maskedcolor,font=\small] at (1.5,2.5) {$\times$};
\filldraw[fill=paper,draw=linegray] (2,2) rectangle (3,3);
\node[maskedcolor,font=\small] at (2.5,2.5) {$\times$};
\filldraw[fill=datacolor!15,draw=linegray] (0,1) rectangle (1,2);
\node[ink,font=\small] at (.5,1.5) {$0.124$};
\filldraw[fill=datacolor!75,draw=linegray] (1,1) rectangle (2,2);
\node[white,font=\small\bfseries] at (1.5,1.5) {$0.876$};
\filldraw[fill=paper,draw=linegray] (2,1) rectangle (3,2);
\node[maskedcolor,font=\small] at (2.5,1.5) {$\times$};
\filldraw[fill=datacolor!20,draw=linegray] (0,0) rectangle (1,1);
\node[ink,font=\small] at (.5,.5) {$0.168$};
\filldraw[fill=datacolor!26,draw=linegray] (1,0) rectangle (2,1);
\node[ink,font=\small] at (1.5,.5) {$0.242$};
\filldraw[fill=datacolor!56,draw=linegray] (2,0) rectangle (3,1);
\node[ink,font=\small] at (2.5,.5) {$0.590$};
\node[anchor=west,align=left,font=\footnotesize] at (3.35,2.3)
 {颜色越深\\读取权重越大};
\node[anchor=west,align=left,font=\footnotesize] at (3.35,.9)
 {$\times$ 为未来位置\\因果 Mask 后权重为 $0$};
\end{tikzpicture}
\captionof{figure}{真实三维主算例的 $A=\operatorname{softmax}(QK^\top/\sqrt3+C)$；每行按可见 Key 归一化}
\label{fig:attention-a-weights}
\end{center}
\tagstructend
